\documentclass[aspectratio=169,11pt]{beamer}

%% Shared preamble for Programming and Quantitative Methods in Finance Beamer presentations
%% Adapted from SPM-PEM/spmcourse-beamer.sty (English localisation)

\usepackage{xcolor}
\usepackage{amsmath,amssymb,bm,mathtools}
\usepackage{booktabs}
\usepackage{hyperref}
\usepackage[english]{babel}

\setbeamertemplate{footline}[frame number]
\setbeamertemplate{itemize items}[circle]
\setbeamercolor{structure}{fg=red!30!green!60!black}
\setbeamercolor{item}{fg=red!30!green!60!black}
\setbeamercolor{itemize item}{fg=red!30!green!60!black}

\definecolor{slidebg0}{RGB}{255,255,255}
\definecolor{slidebg1}{RGB}{220,220,220}
\definecolor{slidebg2}{RGB}{220,235,255}
\definecolor{slidebg3}{RGB}{220,255,230}
\definecolor{slidebg4}{RGB}{255,235,210}
\definecolor{slidebg5}{RGB}{235,220,255}
\definecolor{slidebg6}{RGB}{255,220,205}

\newcounter{slidebg}
\setcounter{slidebg}{1}
\newcommand{\alternateslidebg}{
  \ifnum\value{slidebg}=1
    \setbeamercolor{background canvas}{bg=slidebg1}\setcounter{slidebg}{2}
  \else\ifnum\value{slidebg}=2
    \setbeamercolor{background canvas}{bg=slidebg2}\setcounter{slidebg}{3}
  \else\ifnum\value{slidebg}=3
    \setbeamercolor{background canvas}{bg=slidebg3}\setcounter{slidebg}{4}
  \else\ifnum\value{slidebg}=4
    \setbeamercolor{background canvas}{bg=slidebg4}\setcounter{slidebg}{5}
  \else\ifnum\value{slidebg}=5
    \setbeamercolor{background canvas}{bg=slidebg5}\setcounter{slidebg}{6}
  \else
    \setbeamercolor{background canvas}{bg=slidebg6}\setcounter{slidebg}{1}
  \fi\fi\fi\fi\fi
}

\ExplSyntaxOn
\NewExpandableDocumentCommand{\tocsectionbg}{m}
  { slidebg\int_eval:n { \int_mod:nn { #1 } { 6 } + 1 } }
\ExplSyntaxOff

\setbeamertemplate{section in toc}{%
  \setbeamercolor{section in toc}{fg=black,bg=\tocsectionbg{\inserttocsectionnumber}}%
  \begin{beamercolorbox}[wd=\textwidth,sep=0.7ex,leftskip=1em,rightskip=1em]{section in toc}
    \inserttocsectionnumber.\enspace\inserttocsection
  \end{beamercolorbox}%
  \vspace{0.3ex}%
}
\newcommand{\newsection}[1]{\begin{frame}[plain]
\centering\vfill
{\usebeamercolor[fg]{structure}\LARGE\bfseries \insertsectionnumber\ #1}
\vfill\end{frame}}
\NewDocumentCommand{\bgsection}{s m}{%
  \alternateslidebg
  \section{#2}%
  \IfBooleanF{#1}{\newsection{#2}}%
}
\newcommand{\titleandoutline}{
  \alternateslidebg
  \frame{\titlepage}
  \begin{frame}{Outline}
  \tableofcontents[hideallsubsections]
  \end{frame}
}

\title{Basics, Variables, and Strings}
\subtitle{Lesson 1: getting started with Python}
\author{Programming and Quantitative Methods in Finance}
\date{}

\begin{document}

\titleandoutline

\bgsection{Getting started}

\begin{frame}{Installing Python}
  \begin{itemize}
    \item Install a current Python 3 release from \url{https://www.python.org/}
          or use a managed distribution such as Anaconda.
    \item During installation, make sure Python can be run from the
          command line. Note that on some systems, you may need to use
          \texttt{python3} or \texttt{Py} instead of \texttt{python}.
    \item Verify the installation with \texttt{python3 --version}.
    \item Use the same interpreter consistently for scripts, notebooks, and
          package installation.
  \end{itemize}
\end{frame}

\begin{frame}{Virtual environment}
  \begin{itemize}
    \item Keep course packages isolated from the system Python installation.
    \item From the repository root, create and activate \texttt{.venv}.
  \end{itemize}
  \begin{block}{macOS and Linux}
    \texttt{python3 -m venv .venv}\\
    \texttt{source .venv/bin/activate}\\
    \texttt{python -m pip install --upgrade pip}
  \end{block}
  \begin{block}{Windows PowerShell}
    \texttt{.venv\textbackslash Scripts\textbackslash Activate.ps1}
  \end{block}
\end{frame}

\begin{frame}{Integrated development environments (IDEs)}
  \begin{description}
    \item[Visual Studio Code] Recommended for this course: scripts, notebooks,
          extensions, and an integrated terminal. This is where exams will be given.
    \item[Jupyter] Excellent for exploratory analysis and presenting code,
          output, and plots together.
    \item[PyCharm] A full-featured Python IDE and a valid alternative.
    \item[Online editors] Useful for quick experiments, but less suitable for
          the complete course repository and local data files.
  \end{description}
\end{frame}

\begin{frame}{Visual Studio Code}
  \begin{enumerate}
    \item Install Visual Studio Code from \url{https://code.visualstudio.com/}.
    \item Install Microsoft's \textbf{Python} extension.
    \item Install Microsoft's \textbf{Jupyter} extension for notebook files.
    \item Open the course repository as a folder.
    \item Use ``Python: Select Interpreter'' and choose the course environment.
  \end{enumerate}
\end{frame}

\begin{frame}[fragile]{In-class task: Hello World}
  Now that Python and VS Code are installed, try your first program.
  \begin{block}{hello.py}
\begin{verbatim}
print("Hello World!")
\end{verbatim}
  \end{block}
  \begin{enumerate}
    \item Create a file called \texttt{hello.py} in the course folder.
    \item Run it with VS Code's ``Run Python File in Terminal'', or with
          \texttt{python3 hello.py} in a terminal.
    \item Confirm the output is exactly \texttt{Hello World!}
  \end{enumerate}
\end{frame}

\begin{frame}[fragile]{In-class task: comments and line continuation}
  \begin{enumerate}
    \item Add a comment above the \texttt{print} line explaining what the
          program does, and an inline comment after it.
  \end{enumerate}
  \begin{block}{Comments}
\begin{verbatim}
# Hello World program
print("Hello World!")  # this line displays the text
\end{verbatim}
  \end{block}
  \begin{enumerate}
    \setcounter{enumi}{1}
    \item Rewrite the \texttt{print} call using a backslash to split it over
          three lines, and run it again.
  \end{enumerate}
  \begin{block}{Line continuation}
\begin{verbatim}
print( \
      "Hello World!" \
     )
\end{verbatim}
  \end{block}
\end{frame}

\bgsection{Jupyter notebooks}

\begin{frame}{What is an \texttt{.ipynb} file?}
  \begin{itemize}
    \item A Jupyter notebook stores code cells, Markdown explanations, and
          output in one document.
    \item It is useful for data analysis, plots, and step-by-step demonstrations.
    \item It is mainly an educational and presentation format in this course;
          submitted classroom work may still be requested as a \texttt{.py} file.
    \item Notebook files are JSON documents internally, but should be edited
          through Jupyter or VS Code rather than by hand.
  \end{itemize}
\end{frame}

\begin{frame}{Running a notebook in VS Code}
  \begin{enumerate}
    \item Open the \texttt{.ipynb} file.
    \item Enable or install the Python and Jupyter extensions if prompted.
    \item Select the course \texttt{.venv} as the notebook kernel.
    \item Run cells individually, or run the notebook from the toolbar.
    \item Re-run relevant cells after changing code or input data.
  \end{enumerate}
  \begin{block}{Kernel reminder}
    If several Python installations exist, the first run may ask which one
    should execute the notebook. Choose the course environment.
  \end{block}
\end{frame}

\begin{frame}{Notebook troubleshooting}
  \begin{description}
    \item[No kernel] Select the \texttt{.venv} interpreter in the notebook's
          kernel picker.
    \item[Missing package] Install the package into the selected environment,
          then restart the kernel.
    \item[Stale output] Re-run the cell, or restart and run all cells in order.
    \item[Different results] Check the selected interpreter, package versions,
          working directory, and input data.
  \end{description}
\end{frame}

\begin{frame}{Notebook habits}
  \begin{itemize}
    \item Run cells in a meaningful order and keep the notebook reproducible.
    \item Explain non-obvious calculations in Markdown cells.
    \item Do not treat displayed output as current until its code has been run
          in the selected kernel.
    \item For submitted work, follow the file format and submission instructions
          given in class.
  \end{itemize}
\end{frame}

\bgsection{Programming basics}

\begin{frame}{Programs: input, processing, output}
  \begin{itemize}
    \item Computers cannot read our minds, so we need a way to communicate
          with them: this is what applications (programs, software) do.
    \item The user gives instructions (\textbf{input}); the computer executes
          them; the computer reports back the result (\textbf{output}).
    \item \textbf{Programming} is the process of building the applications
          that carry out this communication.
  \end{itemize}
\end{frame}

\begin{frame}{Machine code, languages, and algorithms}
  \begin{itemize}
    \item Computers understand \textbf{machine code}: instructions and data in
          a form the processor can execute (binary, for today's machines).
    \item Machine code is very hard for humans to read, so we use
          \textbf{programming languages} instead: sets of textual instructions
          that, combined according to a formal structure (\textbf{syntax}),
          produce machine code.
    \item Low-level (\textbf{assembly}) languages are close to machine code and
          hard to read; we use \textbf{high-level} languages (Python, C\#, and
          similar), designed to be easy for humans to read and write.
    \item An \textbf{algorithm} is an abstract, finite, precisely defined
          sequence of steps that solves a problem (for example: find the
          smallest value in a set of numbers).
    \item \textbf{Implementation} is writing that algorithm in a programming
          language so that a working application results.
  \end{itemize}
\end{frame}

\begin{frame}{Python}
  \begin{itemize}
    \item Python is a general-purpose, high-level programming language created
          by Guido van Rossum and first released in 1991.
    \item It is interpreted: code runs line by line, so syntax errors only
          appear once the interpreter reaches them.
    \item Python's simple, readable syntax has made it very widely used,
          including for demanding tasks such as machine learning.
    \item An active community and extensive documentation also mean AI
          assistants have a great deal of good Python code to learn from.
  \end{itemize}
  \begin{block}{A note on AI tools}
    Using AI assistants for programming tasks is encouraged in this course.
    But to handle advanced problems and to direct an AI tool well, you first
    need to master the basics yourself --- and that only happens if you do not
    let a machine do the thinking for you.
  \end{block}
\end{frame}

\bgsection{Variables}

\begin{frame}[fragile]{Python variables}
  \begin{itemize}
    \item A variable is a named container for data: think of it as a labelled
          box you can fill and later retrieve.
    \item Unlike many languages, Python has no keyword to declare a variable.
          A variable is created the first time a value is assigned to it.
  \end{itemize}
  \begin{block}{Example}
\begin{verbatim}
x = 1
x = "Corvinus"
print(x)
\end{verbatim}
    Output: \texttt{Corvinus}
  \end{block}
  \begin{itemize}
    \item Strings may be written with single or double quotes: no difference.
    \item An assignment may reuse other variables, or even the variable itself:
          \texttt{x = x + y + 1}. Self-concatenation can be shortened with
          \texttt{+=}.
  \end{itemize}
\end{frame}

\begin{frame}[fragile]{Variable names}
  \begin{itemize}
    \item A name can be short (\texttt{x}, \texttt{y}) or descriptive
          (\texttt{age}, \texttt{car\_name}, \texttt{total\_weight}).
  \end{itemize}
  \textbf{Rules:}
  \begin{itemize}
    \item Must start with a letter or an underscore \texttt{\_}.
    \item Cannot start with a digit.
    \item May only contain letters, digits, and underscores.
    \item Names are \textbf{case-sensitive}: \texttt{age}, \texttt{Age}, and
          \texttt{AGE} are three different variables.
    \item Python 3 allows many Unicode symbols in names.
    \item A name cannot be a reserved word (an existing keyword, method, or
          class name).
  \end{itemize}
  \begin{block}{A misspelled name causes a runtime \texttt{NameError}}
\begin{verbatim}
message = "Hello World!"
print(mesage)
\end{verbatim}
    \texttt{NameError: name 'mesage' is not defined}
  \end{block}
\end{frame}

\begin{frame}{Naming conventions}
  \begin{itemize}
    \item These are not required by the language, but are good practice.
    \item Global variable names should be all lowercase, with words separated
          by underscores (for now we only use global variables).
    \item Avoid names that are too generic or overly complex: find the
          balance. Bad: \texttt{data\_structure}, \texttt{my\_list},
          \texttt{dictionary\_for\_the\_purpose\_of\_storing\_word\_definitions}.
    \item When using CamelCase, capitalise each letter of an abbreviation
          (e.g.\ \texttt{HTTPServer}).
  \end{itemize}
\end{frame}

\begin{frame}{Python's built-in data types}
  \small
  \begin{tabular}{ll}
    \toprule
    Category & Type name(s) \\
    \midrule
    String (text) & \texttt{str} \\
    Numeric types & \texttt{int}, \texttt{float}, \texttt{complex} \\
    Boolean type & \texttt{bool} \\
    Sequence types & \texttt{list}, \texttt{tuple}, \texttt{range} \\
    Dictionary & \texttt{dict} \\
    Set & \texttt{set}, \texttt{frozenset} \\
    Binary types & \texttt{bytes}, \texttt{bytearray}, \texttt{memoryview} \\
    \bottomrule
  \end{tabular}
\end{frame}

\begin{frame}[fragile]{Checking a variable's type}
  \begin{itemize}
    \item A variable's type never needs to be declared, and it may even change
          after assignment.
    \item The \texttt{type()} function reports a variable's current type.
  \end{itemize}
  \begin{block}{Example}
\begin{verbatim}
x = 1
print(type(x))   # <class 'int'>
print(x)         # 1

x = "Corvinus"
print(type(x))   # <class 'str'>
print(x)         # Corvinus
\end{verbatim}
  \end{block}
\end{frame}

\bgsection{Strings}

\begin{frame}[fragile]{String literals}
  \begin{itemize}
    \item Python strings are surrounded by a pair of single or double quotes.
          A directly written character sequence like this is a \textbf{literal}.
    \item \texttt{"hello"} is exactly the same as \texttt{'hello'}.
    \item Using the other quote character lets you include a quote mark in the
          text: \texttt{print("D'Artagnan is not one of the three musketeers")}.
    \item The closing quote must match the opening one; otherwise Python
          reports a syntax error (the string never closes before end of line).
  \end{itemize}
  \begin{block}{Triple-quoted strings}
    Triple quotes (\texttt{"""..."""}) allow writing a long string across
    several source lines for readability. The displayed text is still a single
    line unless newline characters are included explicitly.
  \end{block}
\end{frame}

\begin{frame}[fragile]{Characters and indexing}
  \begin{itemize}
    \item A string behaves like a list: any element can be accessed by index.
    \item Like most languages, Python indexes strings starting at 0.
    \item Python has no separate character type: a single character is just a
          one-letter string.
  \end{itemize}
  \begin{block}{Example}
\begin{verbatim}
x = "Hello :)"
y = x[0]
print(y)          # H
print(type(y))    # <class 'str'>
\end{verbatim}
  \end{block}
\end{frame}

\begin{frame}[fragile]{Searching within a string}
  \begin{itemize}
    \item The \texttt{in} operator checks whether a substring is present and
          returns a boolean.
    \item \texttt{find()} returns the position of the first match (or \texttt{-1}).
  \end{itemize}
  \begin{block}{Example}
\begin{verbatim}
x = "Budapest Corvinus University"
print("Corvinus" in x)        # True
start = x.find("Corvinus")
print(start)                  # 9
\end{verbatim}
  \end{block}
\end{frame}

\begin{frame}[fragile]{Slicing}
  Slicing uses \texttt{string[start:end:step]}.
  \begin{itemize}
    \item \texttt{start}: index of the first character to keep (default 0).
    \item \texttt{end}: index just past the last character to keep (default:
          the string's length).
    \item \texttt{step}: keep every \texttt{step}-th character (default 1).
  \end{itemize}
  \begin{block}{Example}
\begin{verbatim}
x = "Budapest Corvinus University"
print(x[:])       # the whole string
print(x[:1])      # "B"
print(x[::2])     # every second character
start = x.find("Corvinus")
print(x[start:start+8])   # "Corvinus"
\end{verbatim}
  \end{block}
\end{frame}

\begin{frame}[fragile]{String functions}
  \begin{itemize}
    \item \texttt{len(x)} returns the length of \texttt{x} (works for other
          types too, not only strings).
    \item Type-specific functions are called with a dot after the variable,
          e.g.\ \texttt{x.strip()}; the variable itself is a hidden first
          argument, so the parentheses are often empty.
  \end{itemize}
  \begin{block}{Common string methods}
\begin{verbatim}
x = " Budapest Corvinus University "
print(x.strip())     # remove leading/trailing spaces
print(x.upper())     # BUDAPEST CORVINUS UNIVERSITY
print(x.lower())     # budapest corvinus university
print("kis pista".title())     # Kis Pista
print(x.replace("Budapest", "Vienna"))
print("Budapest".upper())  # methods work on literals too
\end{verbatim}
  \end{block}
\end{frame}

\begin{frame}[fragile]{Concatenation and \texttt{str.format()}}
  \begin{itemize}
    \item The \texttt{+} operator joins strings together.
    \item \texttt{str.format()} is more readable for building longer messages,
          and can insert non-string values (e.g.\ numbers) without manual
          type conversion.
  \end{itemize}
  \begin{block}{Example}
\begin{verbatim}
x, y, z = "Budapest", "Corvinus", "University"
print(x + " " + y + " " + z)

s = "My name is {}. {} {}."
print(s.format("Bond", "James", "Bond"))

s = "My name is {0}. {1} {0}."   # numbered placeholders
print(s.format("Bond", "James"))
\end{verbatim}
  \end{block}
\end{frame}

\begin{frame}[fragile]{f-strings and format parameters}
  \begin{itemize}
    \item Since Python 3.6, an \texttt{f} or \texttt{F} before the opening
          quote creates an \textbf{f-string}: variables can be written
          directly inside \texttt{\{\}}, with no function call needed.
    \item Both approaches accept format parameters, most often used for
          numbers and dates.
  \end{itemize}
  \begin{block}{Example}
\begin{verbatim}
first, last = "James", "Bond"
print(f"My name is {last}. {first} {last}.")

number = 5
print(f"Short decimal: {number:.2f}")   # Short decimal: 5.00
print(f"Long decimal: {number:.10f}")   # Long decimal: 5.0000000000
\end{verbatim}
  \end{block}
\end{frame}

\begin{frame}[fragile]{Escape characters and raw strings}
  \begin{itemize}
    \item The backslash \texttt{\textbackslash} is the \textbf{escape}
          character: it tells Python the next character has a special meaning
          rather than appearing literally.
    \item An escaped quote mark is treated as an ordinary character, not as
          the string's opening or closing mark.
    \item To print a literal backslash, write it twice: \texttt{\textbackslash\textbackslash}.
    \item Common special characters: \texttt{\textbackslash n} (newline) and
          \texttt{\textbackslash t} (tab).
    \item An \texttt{r} or \texttt{R} before the opening quote creates a
          \textbf{raw string}, which disables escape-character handling ---
          useful for file paths full of backslashes.
  \end{itemize}
  \begin{block}{Example}
\begin{verbatim}
print("D\'Artagnan and the \"Three Musketeers\".")
print("c:\\temp\\file.txt")
print(r"c:\temp\file.txt")   # raw string: backslashes kept as-is
\end{verbatim}
  \end{block}
\end{frame}

\bgsection{In-class exercises}

\begin{frame}{In-class exercises: variables}
  \begin{enumerate}
    \item Store your surname and given name in two separate, all-lowercase
          variables, then print your full name with a single \texttt{print}
          call.
    \item Assign an integer to a variable, print its \texttt{type()}, then
          reassign a string to the same variable and print its \texttt{type()}
          again.
    \item Intentionally misspell a variable name in a \texttt{print} call, run
          the program, and read the \texttt{NameError} message that results.
  \end{enumerate}
\end{frame}

\begin{frame}{In-class exercises: strings (1/2)}
  \small
  \begin{enumerate}
    \item You are given a string containing two lines of text separated only
          by a single space, for example: \texttt{"The sun scorches the field
          Grasshoppers graze there"}. Modify the code so that \texttt{print}
          shows the two parts on separate lines.
    \item What does the \texttt{r} prefix before a string literal mean?
    \item Which type is suitable in Python for storing a single character?
    \item Write a program that swaps the two words of a two-word string
          separated by a single space.
    \item Write a program that reverses a string, e.g.\ \texttt{ABC} becomes
          \texttt{CBA}. Hint: the step in a slice can be negative.
  \end{enumerate}
\end{frame}

\begin{frame}{In-class exercises: strings (2/2)}
  \small
  \begin{enumerate}
    \setcounter{enumi}{5}
    \item Write a program that upper-cases the first word of a two-word name,
          e.g.\ \texttt{Jane Doe} $\rightarrow$ \texttt{JANE Doe}.
    \item Write a program that extracts the initials from a two-word name,
          e.g.\ \texttt{Jane Doe} $\rightarrow$ \texttt{JD}.
    \item Write a program that fixes the capitalisation of a lowercase
          two-word name, e.g.\ \texttt{jane doe} $\rightarrow$ \texttt{Jane Doe}.
          Try it both with and without the dedicated string method.
    \item Write a program that strips the leading and trailing spaces from a
          string and then prints the length of the result, using as few lines
          of code as possible. Try it both with and without the dedicated
          string method.
  \end{enumerate}
\end{frame}

\end{document}
